\newpage
\phantomsection
\label{prf:riemann-integrable-implies-hk-integrable}
\LRAProofFor{thm:riemann-integrable-implies-hk-integrable}
\begin{remark*}[Return]\hyperref[thm:riemann-integrable-implies-hk-integrable]{Return to Theorem}\end{remark*}
\begin{theorem*}[Riemann Integrable Implies Henstock--Kurzweil Integrable]
If \(f\) is Riemann-integrable on \([a,b]\), then \(f\) is
Henstock--Kurzweil integrable on \([a,b]\), with the same value.
\end{theorem*}
\begin{proof}[Professional Standard Proof]
\LRAProofBodyStart
Let \(L\) be the Riemann integral of \(f\). Given \(\varepsilon>0\), choose
\(\eta>0\) so that every tagged partition with mesh less than \(\eta\) has
Riemann sum within \(\varepsilon\) of \(L\). Use the constant gauge
\(\delta(x)=\eta/2\). Every \(\delta\)-fine tagged partition then has mesh
less than \(\eta\), so its tagged sum is within \(\varepsilon\) of \(L\).
\end{proof}
\begin{proof}[Detailed Learning Proof]
\LRAProofBodyStart
The Henstock--Kurzweil definition allows more gauges than Riemann needs. To
recover Riemann's situation, take a gauge that is constant at every point. A
\(\delta\)-fine partition for that constant gauge has every subinterval shorter
than the Riemann mesh tolerance, so the old Riemann estimate applies without
change.
\end{proof}
\begin{remark*}[Proof structure]
Embed the uniform mesh condition into the gauge condition by choosing a
constant gauge.
\end{remark*}
\begin{dependencies}
\begin{itemize}
  \item \hyperref[def:riemann-integral]{Riemann Integral}.
  \item \hyperref[def:gauge-on-interval]{Gauge on an Interval}.
  \item \hyperref[def:delta-fine-tagged-partition]{Delta-Fine Tagged Partition}.
  \item \hyperref[def:henstock-kurzweil-integral]{Henstock--Kurzweil Integral}.
\end{itemize}
\end{dependencies}
\clearpage
